\exercisegroup

\begin{enumerate}
\def\labelenumi{\arabic{enumi})}
\tightlist
\item
  \begin{enumerate}
  \def\labelenumii{\alph{enumii})}
  \tightlist
  \item
    设 E 是向量空间。我们说 E 中一个（以 0 为顶点的）尖凸锥 C 是极大的，如果 C 是按包含排序的、以 0 为顶点且 \(\neq E\) 的尖凸锥所成之集中的一个极大元。证明：尖凸锥 C 是极大的，其充要条件是它是由一个过 0 的超平面定义的闭半空间。为建立这一结果，相继证明极大尖凸锥 C 的下列性质；
  \end{enumerate}
\end{enumerate}

α) 我们有 C ∪ (− C) = E（用反证法论证）。

\(\beta)\) 若 \(z\) 是 C 的一个非代数内点（II, 第 26 页），则 \(-z \in C\)（同法）。由此推出 C 含有代数内点。

\(\gamma)\) 含于 C 中的最大向量子空间 \(H = C \cap (-C)\) 是一个超平面。（过渡到商空间 F = E/H，这就化归为利用 \(\beta)\) 证明：若 C 的除顶点外的所有点都是代数内点，则 E 必为 1 维。）

\begin{enumerate}
\def\labelenumi{\alph{enumi})}
\setcounter{enumi}{1}
\tightlist
\item
  给出一个极大非尖凸锥（在以 0 为顶点的非尖凸锥之集中）的例子，它没有代数内点（参见 II, 第 65 页, 习题 5）。
\end{enumerate}

\begin{enumerate}
\def\labelenumi{\arabic{enumi})}
\setcounter{enumi}{1}
\item
  设 N 是向量空间 E 的向量子空间 M 中的一个超平面，A 是 E 中的凸集，使 \(A \cap M\) 的所有点都位于 N 的同一侧，且它还具有下列性质：对 E 中任何 \(y \neq 0\)，存在 \(x \in A \cap M\)，使 \(x + \lambda y \in A\) 对一切 \(\lambda\)（\(|\lambda|\) 充分小的）成立。证明：此时存在 E 的一个超平面 H，使 A 的所有点都位于 H 的同一侧，且 \(H \cap M = N\)。（化归到 \(N = \{0\}\) 的情形；若 \(a \neq 0\) 属于 \(A \cap M\)，考察含 A 且不含 -a 的以 0 为顶点的尖凸锥所成之集 U；证明 U 存在极大元 C，且 C 是极大尖凸锥（习题 1）。）由此给出 Hahn-Banach 定理的一个新证明。
\item
  设 A 是拓扑向量空间 E 中的凸集，\(x_{0}\) 是 E 的一点。那么，存在一个闭超平面 H，包含 \(x_{0}\)，且使 A 的所有点都位于 H 的同一侧，其充要条件是：存在一个以 \(x_{0}\) 为顶点的非尖凸锥 C，它至少含一个内点且不与 A 相交。（关于不含于任何由超平面定义的半空间中的凸集 \(A \neq E\) 的例子，见 II, 第 65 页, 习题 5。）
\item
  设 E 是赋范空间，A 是对由 E 诱导的一致结构完备的凸集。
\end{enumerate}

\begin{enumerate}
\def\labelenumi{\alph{enumi})}
\item
  设 \(x'\) 是 \(E\) 上的连续线性型，它在 \(A\) 中有界。考察一个数 \(k > 0\) 以及一个闭凸锥 \(P\)（\(E\) 中以 0 为顶点、由 \(x \in E\)（满足 \(\| x \| \leqslant k < \langle x, x' \rangle\)）组成的）；它是尖的且是真的。证明：对 \(E\) 上以 \(P\) 为诸元素 \(\geqslant 0\) 所成之集的那个序，集 \(A\) 是归纳的（利用 \(x'\) 在 \(A\) 上的限制递增且有界这一事实）。
\item
  由 a) 推出：A 的边界 F 中属于 A 的某个支承超平面的点所成之集在 F 中稠密（Bishop-Phelps 定理）。（对每一点 \(z \in F\)，考察一点 \(y \in \complement A\)，它与 z 任意接近，并用一个方程为 \(\langle x, x' \rangle = \alpha\) 的闭超平面把 y 与 A 严格分离，其中 \(\|x' \| = 1\)，然后取 k > 1 用 a)，并用上面的习题 3。）
\end{enumerate}

\begin{enumerate}
\def\labelenumi{\arabic{enumi})}
\setcounter{enumi}{4}
\tightlist
\item
  设 \(\mathbf{A}\) 是 \(\mathbf{R}^n\) 中一个闭凸集，\(x_0\) 是 \(\complement \mathbf{A}\) 的一点；把 \(\mathbf{R}^n\) 中的欧氏距离记作 \(d\)。
\end{enumerate}

\begin{enumerate}
\def\labelenumi{\alph{enumi})}
\item
  不用 II, 第 36 页的 th. 1，证明：存在一点且仅一点 \(x \in \mathbf{A}\)，使 \(d(x_0, x) = d(x_0, \mathbf{A})\)，且与连接 \(x_0\) 与 \(x\) 的直线正交且过 \(x\) 的超平面是 \(\mathbf{A}\) 的一个支承超平面。
\item
  当空间 E 为有限维时，由 a) 给出 II, 第 36 页 th. 1 的一个新证明。（化归到 M 为 A 的边界点 \(x_{0}\) 的情形；注意 \(x_{0}\) 与 A 的诸支承超平面的距离的下确界为零，并利用 \(S_{n-1}\) 的紧性。）
\end{enumerate}

\begin{enumerate}
\def\labelenumi{\arabic{enumi})}
\setcounter{enumi}{5}
\item
  设 A 是 \(R^{n}\) 中的闭集，具有下列性质：对每个 \(x \in R^{n}\)，存在一点且仅一点 \(y \in A\)，使 \(d(x, y) = d(x, A)\)，其中 d 是欧氏距离。证明：A 是凸的。（用反证法，考察 A 中以 a、b 为端点且含一点 \(c \in \complement A\) 的闭线段；存在一个以 c 为中心、含于 \(\complement A\) 的闭球 B；考察包含 B 且其内部不与 A 相交的闭球 S 所成之集；证明这些球的半径上有界，并由此推出其中存在一个球 \(S_{0}\)，其半径 \(\rho\) 尽可能大。然后证明 \(S_{0}\) 只能与 A 交于一点，而这蕴含 B 中存在一个半径 \(> \rho\) 的球，由此得出矛盾。）
\item
  在 Hausdorff 局部凸空间 \(\mathbf{E}\) 中，设 \(\mathbf{A}\) 是完备凸集，\(\mathbf{B}\) 是预紧闭凸集，满足 \(\mathbf{A} \cap \mathbf{B} = \varnothing\)。证明：存在一个把 A 与 B 分离的闭超平面（在完备化 \(\hat{\mathbf{E}}\) 中论证）。考察 A 为有限维的情形。
\item
  在 Hausdorff 局部凸空间中，设 A 与 B 是两个无公共点的闭凸集，满足 \(C_{A} \cap C_{B} = \{0\}\)（II, 第 67 页, 习题 14），且 B 局部紧。证明：存在一个把 A 与 B 分离的闭超平面（参见 II, 第 67 页, 习题 16）。类似地，若 A 与 B 是两个以 0 为顶点的闭凸锥，满足 \(A \cap B = \{0\}\) 且 B 局部紧，则存在一个过 0 且把 A 与 B 分离的闭超平面（利用 II, 第 39 页的引理 1）。
\item
  由习题 8 推出：若 V 是 E 的有限维向量子空间，C 是 E 中以 0 为顶点的闭凸锥，满足 \(C \cap V = \{0\}\)，则存在 C 的一个包含 V 的支承超平面（利用 II, 第 39 页的引理 1）。
\item
  在赋范空间 \(E = l^{1}(\mathbf{N})\)（实数 \(x = (\xi_{n})_{n \in \mathbf{N}}\) 的可和序列所成的）中，设 D 是由关系 \(\xi_{n} = 0\)（对 \(n \geqslant 1\)）定义的直线。证明：存在实数 > 0 的两个递增序列 \((\alpha_{n})\)、\((\beta_{n})\)，使得由不等式 \(\xi_{0} \geqslant |\alpha_{n}\xi_{n} - \beta_{n}|\)（对 \(n \geqslant 1\)）定义的凸集 A 是闭的、无界的、不与 D 相交，且不存在把 A 与 D 分离的闭超平面（选取 \(\alpha_{n}\) 与 \(\beta_{n}\) 使 A - D 处处稠密）。
\end{enumerate}

* 11) a) 设 E 是 Hilbert 空间，F 是 E 的对偶 E' 的一个异于 E' 的处处稠密子空间；E 的单位球 B 对弱拓扑 σ(E, F) 是紧的，且单位球面上存在一点 a，过它没有 B 的（在 σ(E, F) 中）闭的支承超平面。

\begin{enumerate}
\def\labelenumi{\alph{enumi})}
\setcounter{enumi}{1}
\item
  给 E 赋以拓扑 \(\sigma(\mathrm{E},\mathrm{F})\)，并在乘积空间 \(\mathbf{G} = \mathbf{E} \times \mathbf{R}\) 中考察点对 \((x,\zeta)\)（满足 \(\|x\| < 1, \zeta \geqslant \|x\| / (1 - \|x\|)\)）所成之集 A。证明：A 是闭的且局部紧的，但若 D 是 G 中方程为 \(x = a\) 的直线，则 \(\mathrm{D} \cap \mathrm{A} = \emptyset\)，且 G 中不存在把 A 与 D 分离的闭超平面。
\item
  证明：当我们给 E 赋以拓扑 \(\sigma(E, F)\) 时，E 的子空间 B 中存在一个连续仿射实值函数，它不是 E 中连续仿射函数在 B 上的限制。 *
\end{enumerate}

\begin{enumerate}
\def\labelenumi{\arabic{enumi})}
\setcounter{enumi}{11}
\item
  在 \(\mathbf{R}^3\) 中考察闭凸锥 \(C\)（由关系 \(\xi_1 \geqslant 0\)、\(\xi_2 \geqslant 0\)、\(\xi_3^2 \leqslant \xi_1\xi_2\) 定义的）。证明：直线 \(D\)（方程为 \(\xi_1 = 0\)、\(\xi_3 = 1\)）不与 \(C\) 相交，但不存在过原点 \(0\)、包含 \(D\) 且不与 \(C - \{0\}\) 相交的平面。
\item
  设 A 与 B 是空间 \(R^{n}\) 中两个闭凸集，使得：若 V 与 W 分别是由 A 与 B 生成的仿射线性流形，则 \(A \cap B\) 中没有一点既是 A 相对于 V 的内点，又是 B 相对于 W 的内点。证明：存在一个把 A 与 B 分离的超平面。（通过取商，化归到下列情形：流形 V、W 之一含于另一个，或者 V 与 W 是 E 中互补的向量子空间。）
\item
  设 A 是不含直线的抛物闭凸集（II, 第 67 页, 习题 17）。证明：若 B 是不与 A 相交的闭凸集，则 \(\mathbf{R}^n\) 中存在一个把 A 与 B 严格分离的超平面（若 \(d\) 是欧氏距离，证明 \(d(\mathrm{A},\mathrm{B}) > 0\)）（参见习题 12。）
\item
  设 S、T 是 \(R^{n}\) 中两个无公共点的有限集，且满足 \(\operatorname{Card}(S \cup T) \geqslant n + 2\)。为使存在一个把 S 与 T 严格分离的超平面，其充要条件是：对每个有限集 \(F \subset S \cup T\)（由 \(n + 2\) 个点组成的），存在一个把 \(F \cap S\) 与 \(F \cap T\) 严格分离的超平面（利用 Helly 定理（II, 第 68 页, 习题 21））。证明：在这个命题中不能把数 \(n + 2\) 换成 \(n + 1\)，且该命题不推广到 S 与 T 为无穷的情形。
\item
  设 A 是 \(R^{n}\) 中有内点的紧集。证明：若 A 的每个边界点都位于 A 的至少一个支承超平面上，则 A 是凸的。（得出矛盾：证明若 x 与 y 是 A 中两点，使以 x、y 为端点的线段不含于 A，而 z 是 A 的不位于该线段上的内点，则在以 x、y、z 为顶点的三角形中存在 A 的一个异于 x 与 y 的边界点。）
\item
  在 \(R^{n}\) 中，设 A 是对称凸集，0 是它的内点，且其边界不含任何真线段。设 H 是一个齐次超平面，D 是与 H 互补的直线。证明：存在一点 \(a \in H \cap A\)，使在 a 处有 A 的一个平行于 D 的支承超平面。
\item
  在拓扑向量空间 \(\mathbf{E}\) 中，设 \(\mathrm{A}_i\)（\(1 \leqslant i \leqslant n\)）是 \(n\) 个非空开凸集。a) 证明：若诸 \(\mathrm{A}_i\) 之并异于 \(\mathbf{E}\)，则每个点 \(x \in \mathbf{E}\)（不属于任何 \(\mathrm{A}_i\) 的）都属于一个余维数为 \(n\) 的闭线性流形（包含 \(x\) 且不与任何 \(\mathrm{A}_i\) 相交的）（对 \(n\) 作归纳论证）。
\end{enumerate}

\begin{enumerate}
\def\labelenumi{\alph{enumi})}
\setcounter{enumi}{1}
\tightlist
\item
  若诸 \(\mathbf{A}_i\) 的交为空，证明：E 中存在一个余维数为 \(n - 1\) 的闭线性流形，它不与任何 \(\mathbf{A}_i\) 相交（同法）。
\end{enumerate}

\begin{enumerate}
\def\labelenumi{\arabic{enumi})}
\setcounter{enumi}{18}
\tightlist
\item
  设 C、C' 是 Hausdorff 拓扑向量空间 E 中两个被一个闭超平面 H 严格分离的闭凸集。设 H' 是同时支承 C 与 C' 的闭超平面，使 C 与 C' 位于 H' 的同一侧。证明：H' 是具有这些性质且包含 H ∩ H' 的唯一超平面，且 H ∩ H' 是 C ∪ C' 的凸包在 H 上的迹 P 的支承超平面。反之，若 C 与 C' 是紧的，则对 P 在 H 中的每个支承超平面 D，存在一个同时支承 C 与 C'、包含 D 且使 C 与 C' 位于同一侧的超平面 H'。
\end{enumerate}

¶ 20) 设 A、B 是 Hausdorff 局部凸空间 E 中两个不相交的闭凸集，H 是把 A 与 B 分离的闭超平面；设 \(A \cap H \neq \varnothing\)，且 \(A \cap H\) 与每条直线的交是紧的。证明：若 A 或 B 局部紧，则 E 中存在 0 的一个邻域 V，使 \((A + V) \cap B\) 为空。（按 A 局部紧还是 B 局部紧分两种情形；在第一种情形，注意存在一个平行于 H 的超平面 \(H'\)，使若 S 是 H 与 \(H'\) 之间的点所成之集，则 \(A \cap S\) 是紧的。在第二种情形，例如设 \(0 \in B \cap H\)；对 E 中 0 的每个邻域 V，考察集 \((A + V) \cap B\)，并相继考察该集对至少一个 V 相对紧的情形与不然的情形，如 II, 第 67 页, 习题 16 那样。）

¶ 21) a) 在 \(\mathbf{R}^n\) 中，设 \(a_i (1 \leqslant i \leqslant n + 1)\) 是 \(n + 1\) 个仿射无关的点。把诸 \(a_i\) 的凸包记作 \(\mathbf{S}\)，把诸 \(a_i\)（\(i \neq k\) 的）的凸包记作 \(\mathrm{F}_k\)（\(1 \leqslant k \leqslant n + 1\)）。对每个 \(k\)，设 \(\mathrm{C}_k\) 是包含 \(\mathrm{F}_k\) 的紧凸集，并假设

S 含于诸 \(C_k\) 之并；证明此时 \(\bigcap_{k=1}^{n+1} C_k \neq \emptyset\)。（用归谬法并

对 \(n\) 作归纳来论证，考察交集 \(\mathbf{C}_{n+1}'\)（诸 \(\mathbf{C}_i\)（指标 \(i \leqslant n\)）的），并设 \(\mathbf{C}_{n+1} \cap \mathbf{C}_{n+1}' = \emptyset\)，这将允许用一个超平面把这两个凸集严格分离。）

\begin{enumerate}
\def\labelenumi{\alph{enumi})}
\setcounter{enumi}{1}
\tightlist
\item
  设 X 是 Hausdorff 拓扑向量空间 E 中的紧凸集，\((\mathrm{C}_{\lambda})_{\lambda\in\mathrm{L}}\) 是含于 X 中的紧凸集的一个族，使得：对每个具有 n（相应地 m）个元素的集 H ⊂ L，诸 \(C_{\lambda}\)（指标 \(\lambda\in H\)）的交（相应地并）非空（相应地等于 X）。证明：若 \(m\leqslant n+1\)，则交 \(\bigcap_{\lambda\in L}C_{\lambda}\) 非空。（这实际上
\end{enumerate}

就是证明：对 L 的任何由 \(p \geqslant m\) 个指标组成的有限集 H，我们有 \(\bigcap_{\lambda \in H} C_{\lambda} \neq \varnothing\)。对

\(p\) 作归纳论证，假设结果已对 \(p - 1\) 个指标证明。然后再用反证法论证，对每个指标 \(i \in \mathbf{H}\) 考察一点 \(a_i \in \bigcap_{\lambda \in \mathbf{H} - \{i\}} \mathrm{C}_\lambda\)，并借助

Helly 定理（II, 第 68 页, 习题 21）证明诸 \(a_{i}\) 生成一个维数为 \(p - 1\) 的线性流形，最后应用 \(a\)。）

¶ 22) 在 \(\mathbf{R}^n\) 中，设 \((\mathbf{C}_i)_{1 \leqslant i \leqslant m}\) 是以 0 为顶点的闭凸锥的一个有限族，使得其中任何 \(n\) 个之和异于 \(\mathbf{R}^n\)。证明：存在一个过 0 的超平面 H，使得对任何指标 \(i\)，\(\mathbf{C}_i\) 中没有一对点被 H 严格分离。（区分两种情形：

\(\alpha)\) 或者存在一个数 \(r < n\) 与 \(r\) 个指标，例如 1, 2, \ldots, \(r\)，使诸 \(C_i\)（\(i \leqslant r\) 的）生成一个锥，它包含一个维数 \(\geqslant r\) 的向量子空间 V。对 \(n\) 作归纳论证，在 V 的正交补上投影。

β) 或者，对所有 r < n，任何 r 个 \(C_{i}\) 生成一个锥 C，使含于 C 中的最大向量子空间 \(C \cap (-C)\) 的维数 < r。然后考察诸锥 \(C_{i}\) 的一个集，它对于含于一个半空间而言是极大的；一个极大集中至少有 n 个锥。设 \(\Gamma\) 是由属于该极大集的诸锥之并生成的锥。若 \(C_{j}\) 是一个不属于所考察的极大集的锥，证明 \(C_{i} \subset -\Gamma\)。为此，用反证法论证，证明在相反的情形，存在 \(-\Gamma\) 的一个边界点（相对于由 \(\Gamma\) 生成的向量子空间），它是 \(C_{j}\) 的内点（相对于由 \(C_{j}\) 生成的向量子空间）。把这样一个点写成个数最少的 s 个向量之和，其中每个向量属于某个锥 \(-C_{i}\)，这些 \(C_{i}\) 是用于定义 \(\Gamma\) 的那一些；于是 \(s \leqslant n - 1\)。最后证明这些锥与 \(C_{j}\) 生成一个包含 \(s + 1\) 维向量子空间的凸锥，与假设矛盾；为此利用 II, 第 65 页的习题 4。）

\begin{enumerate}
\def\labelenumi{\arabic{enumi})}
\setcounter{enumi}{22}
\item
  设 \(\mathbf{E}\) 是拓扑向量空间，\(\mathcal{T}\) 是 \(\mathbf{E}\) 上的局部凸拓扑，它是粗于给定拓扑 \(\mathcal{T}_0\)（\(\mathbf{E}\) 上的）的一切拓扑中的最细者。若 \(\mathbf{F}\) 是局部凸空间，则 \(\mathbf{E}\) 到 \(\mathbf{F}\) 中的连续线性映射对 \(\mathcal{T}_0\) 与对 \(\mathcal{T}\) 是相同的。\(\mathbf{E}\) 上存在异于零形式的连续线性型，其充要条件是：存在 0 的一个对 \(\mathcal{T}_0\) 而言的邻域，其凸包（对 \(\mathcal{T}_0\)）不是处处稠密的（参见 I, 第 25 页, 习题 4）。
\item
  设 \(\mathbf{E}\) 是无穷维可度量化局部凸空间。
\end{enumerate}

\begin{enumerate}
\def\labelenumi{\alph{enumi})}
\item
  证明：存在 E 的点的一个趋于 0 的序列 \((a_{n})\)，以及 E 的闭向量子空间的一个递减序列 \((\mathbf{L}_n)\)，使 \(\mathbf{L}_n\) 在 E 中的余维数为 \(n\)，且对所有 \(n\)，点 \(a_{n}\) 属于 \(\mathbf{L}_n - \mathbf{L}_{n + 1}\)。
\item
  进一步假设 E 完备。证明：此时可以找到满足条件 a) 的序列 \((a_{n})\) 与 \((\mathrm{L}_{n})\)，使此外对实数的每个有界序列 \((\lambda_{n})\)，通项为 \(\lambda_{n}a_{n}\) 的级数在 E 中交换收敛，且线性映射 \((\xi_{n})\mapsto\sum_{n}\xi_{n}a_{n}\)（Banach 空间 \(\mathcal{B}(\mathbf{N})\) 到 E 中的）是单射且连续的。
\item
  由 b) 推出：当 E 是无穷维 Fréchet 空间时，E 在 R 上的每个基的基数至少等于 \(2^{\mathrm{Card}(\mathbf{N})}\)（参见 I, 第 22 页, 习题 5）。
\end{enumerate}

若存在一个在 \(\mathbf{E}\) 中稠密的可数集，则 \(\mathbf{E}\) 的每个基都具有连续统的基数。

\begin{enumerate}
\def\labelenumi{\arabic{enumi})}
\setcounter{enumi}{24}
\item
  设 E 是可数型的无穷维 Fréchet 空间（因而具有可数的处处稠密子集）（参见 I, 第 25 页, 习题 1）。证明：存在 E 的一个处处稠密超平面 H，它与 E 的每个闭无穷维线性流形相交。（利用这些流形的每个方向子空间中具有连续统基数的基的存在性（习题 24, c)），以及 E 的闭无穷维线性流形所成之集也具有连续统基数这一事实（GT, IX, § 5, 习题 17）；然后应用由 S, III, § 6, 习题 24 得出的 E 上线性型的构造方法。）超平面 H 不含任何无穷维闭向量子空间。
\item
  设 \(\mathbf{E}\) 是可数型的无穷维 Fréchet 空间。
\end{enumerate}

\begin{enumerate}
\def\labelenumi{\alph{enumi})}
\item
  证明：存在 E 的线性无关元素的一个序列 \((a_{n})\)，使序列 \((a_{2n})\) 与 \((a_{2n + 1})\) 中的每一个都是完全的（利用习题 24, c)）。
\item
  设 \(\mathbf{F}\) 是 \(\mathbf{E}\) 的由诸 \(a_{2n + 1}\)（\(n \in \mathbf{N}\)）生成的向量子空间。对每个 \(n > 0\)，设 \(\mathbf{M}_n\) 是由诸 \(a_{2k}\)（\(k \leqslant n\)）生成的子空间。对每个 \(n\)，设 \(\phi_n\) 是 \(\mathbf{F}\) 上的如下映射：\(\mathbf{E}\) 到 \(\mathrm{E / M}_n\) 上的典范同态的限制，并设 \(\mathcal{T}_n\) 是 \(\mathbf{F}\) 上的拓扑，即 \(\phi_n\) 下 \(\mathrm{E / M}_n\) 上商拓扑的逆像。证明：诸拓扑 \(\mathcal{T}_n\) 中的每一个都是 \(\mathbf{F}\) 上的 Hausdorff 局部凸拓扑，但诸 \(\mathcal{T}_n\) 在 \(\mathbf{F}\) 上局部凸拓扑之集中的下确界是 \(\mathbf{F}\) 上最粗的拓扑。
\end{enumerate}

* c) 取 E 为 Hilbert 空间；证明：可以选取序列 \((a_{n})\)，使若 G 是由诸 \(a_{4n + 1}\) 生成的闭向量子空间，则 G 有无穷余维数，且诸 \(a_{2n}\) 与诸 \(a_{4n + 3}\) 在 E/G 中的像仍线性无关。用 \(\mathbf{G}_n\) 表示 E 的子空间，即 G 与由诸 \(a_{4k + 3}\)（\(k\leqslant n\)）生成的子空间之和，并给 \(G_{n}\) 赋以 \(G_{n}\) 上限制的典范映射下 \(E/M_{n}\) 上商拓扑的逆像拓扑。证明：序列 \((\mathbf{G}_{n})\) 是拓扑向量空间的一个归纳系，使 \(G_{n}\) 对 \(G_{n+1}\) 的拓扑在 \(G_{n+1}\) 中是闭的，但 \(G_{n}\) 在该序列的归纳极限空间中不是闭的。 *

\begin{enumerate}
\def\labelenumi{\arabic{enumi})}
\setcounter{enumi}{26}
\tightlist
\item
  设 E、F 是两个 Hausdorff 拓扑向量空间，X（相应地 Y）是 E（相应地 F）中的紧凸集。设 f 是定义在 \(X \times Y\) 中的实值函数，具有下列性质：
\end{enumerate}

\begin{enumerate}
\def\labelenumi{(\roman{enumi})}
\item
  对所有 \(x \in \mathbf{X}\)，映射 \(y \mapsto f(x, y)\) 在 \(\mathbf{Y}\) 中下半连续，且对所有 \(c \in \mathbf{R}\)，由 \(y \in \mathbf{Y}\)（满足 \(f(x, y) \leqslant c\) 的）所成之集是凸的。
\item
  对所有 \(y \in \mathbf{Y}\)，映射 \(x \mapsto f(x, y)\) 在 \(\mathbf{X}\) 中上半连续，且对所有 \(c \in \mathbf{R}\)，由 \(x \in \mathbf{X}\)（满足 \(f(x, y) \geqslant c\) 的）所成之集是凸的。
\end{enumerate}

证明：在这些条件下，我们有

\[
\sup _ {x \in X} (\inf _ {y \in Y} f (x, y)) = \inf _ {y \in Y} (\sup _ {x \in X} f (x, y)).
\]

（用反证法论证，假设存在一个数 \(c\)，使得

\[
\sup _ {x \in \mathbf {X}} (\inf _ {y \in \mathbf {Y}} f (x, y)) <   c <   \inf _ {y \in \mathbf {Y}} (\sup _ {x \in \mathbf {X}} f (x, y)).
\]

对所有 \(x \in X\)（相应地所有 \(y \in Y\)），设 \(A_{x}\) 是由 \(y \in Y\)（满足 \(f(x, y) > c\) 的）所成之集（相应地设 \(B_{y}\) 是由 \(x \in X\)（满足 \(f(x, y) < c\) 的）所成之集），它在 Y 中（相应地在 X 中）是开的；当 x 在 X 中变化（相应地 y 在 Y 中变化）时，诸 \(A_{x}\)（相应地诸 \(B_{y}\)）构成 Y（相应地 X）的一个覆盖。证明：存在两个有限集 \(X_{0} \subset X\)、\(Y_{0} \subset Y\)，使：\(1^{\circ}\) 对属于凸包 \(B_{0}\)（\(Y_{0}\) 的）的每个 y，存在 \(x \in X_{0}\)，使 \(f(x, y) > c\)，且 \(X_{0}\) 对该性质是最小的；\(2^{\circ}\) 对属于凸包 \(A_{0}\)（\(X_{0}\) 的）的每个 x，存在 \(y \in Y_{0}\)，使 \(f(x, y) < c\)，且 \(Y_{0}\) 对该性质是最小的。然后，对所有 \(y \in Y_{0}\)，设 \(C_{y}\) 是由 \(x \in A_{0}\)（满足 \(f(x, y) \geqslant c\) 的）所成之集；利用 II, 第 79 页, 习题 21, a) 证明：诸 \(C_{y}\)（\(y \in Y_{0}\) 的）之交非空。在 \(B_{0}\) 中同样进行，并得出矛盾。）

¶ 28) 设 X 是 Hausdorff 局部凸空间 E 的紧凸子集，f 是 X 中的上半连续凸函数。证明：X 中的连续凸函数 g（满足 \(g(x) > f(x)\)（对所有 \(x \in X\)））所成之集 L 是向下有向的，且其下包络等于 f。（设 u、v 是 L 的元素。为构造 L 中一个小于 u 与 v 的元素，利用与 II, 第 40 页, prop. 6 类似的推理。把该论证中类似于集 K 的集 \(K_{1}\) 解释为位于一个下半连续函数的图象上方的点所成之集，该函数小于 u 与 v 且在每一点处严格大于 f；对这个函数应用 II, 第 39 页的 prop. 5 与 Dini 定理。为证 L 的下包络是 f，注意 f 被一个常数 b 上有界；设 \((x, t)\) 是 E × R 中位于 f 的图象上方的一点，设 \(K'\) 是 \(\{(x, t)\} \cup (X' \times \{b\})\) 的凸包，其中 \(X'\) 是 x 在 X 中的一个适当的紧邻域；对 \(K'\) 像上面对 \(K_{1}\) 那样论证。）

\begin{enumerate}
\def\labelenumi{\arabic{enumi})}
\setcounter{enumi}{28}
\item
  设 X 是 Hausdorff 局部凸空间 E 中的紧凸集。设 u 是 X 中的下半连续凸函数，v 是 X 中的上半连续凹函数，使 \(u(x) > v(x)\) 对所有 \(x \in X\) 成立。则存在一个在 E 中连续的仿射线性函数 f，使 \(v(x) < f(x) < u(x)\) 对所有 \(x \in X\) 成立。
\item
  设 \(\mathbf{X}\) 是 Hausdorff 局部凸空间 \(\mathbf{E}\) 中的紧凸集。证明：\(\mathbf{X}\) 中的下半连续凸函数所成之集是一个格。
\end{enumerate}
% semantic-fragments: legacy-input-seam
\relax{}
